% !TeX root = courtade-kumar.tex
\section{Shannon and Tsallis information}
\label{sec:quadi}
\iffalse
We develop the quantities from \autoref{sec:preliminaries}
for random variables on an arbitrary finite alphabet $\Omega$,
and prove \autoref{lem:tsallis-properties} below.

\subsection*{Divergence and entropy}
For random variables $X,W$ on $\Omega$, define the quadratic
divergence and the quadratic entropy by
\begin{equation*}
  \tsdiv{X}{W}
  =2\sum_{\omega\in\Omega}\nparen{\Pr[X=\omega]-\Pr[W=\omega]}^2,
  \qquad
  \tsent{X}=2\nparen{1-\sum_{\omega\in\Omega}\Pr[X=\omega]^2}.
\end{equation*}
Thus $\tsdiv{X}{W}\ge0$ with equality iff $X$ and $W$ have the
same law, and $\tsent{X}=2\Pr[X\ne X']$ for an independent copy
$X'$ of $X$; up to the factor $2$, $\tsent{\cdot}$ is the Tsallis
$2$-entropy \cite{Tsallis1988}. If $W$ is uniform on $\Omega$,
then
\begin{equation*}
  \tsent{X}=2\nparen{1-\tfrac1{|\Omega|}}-\tsdiv{X}{W}.
\end{equation*}
For bits with $\Pr[X=1]=a$ and $\Pr[W=1]=b$, we have
$\tsdiv{X}{W}=4(a-b)^2$, hence
\begin{equation*}
  \tsdivb{a}{b}=4(a-b)^2,
  \qquad
  \tsent{a}=4a(1-a)=1-\tsdivb{a}{1/2},
\end{equation*}
recovering the definitions of \autoref{sec:preliminaries}.

Conditioning on a random variable $Z$ means averaging over it:
\begin{equation*}
  \tsent{X\emid Z}=\E_{z\sim Z}\tsent{X\emid Z=z},
  \qquad
  \tsdiv{X\emid Z}{W}=\E_{z\sim Z}\tsdiv{X\emid Z=z}{W}.
\end{equation*}

\subsection*{Information}
Let $X,W$ take values in $\Omega$ and $Z$ be an observed random
variable. The key fact is a Pythagorean identity.
\begin{lemma}
\label{lem:quad-pythagoras}
\begin{equation*}
  \tsdiv{X\emid Z}{W}=\tsdiv{X\emid Z}{X}+\tsdiv{X}{W}.
\end{equation*}
\end{lemma}
\begin{proof}
For each $\omega$, write $\Pr[X=\omega\emid Z=z]-\Pr[W=\omega]$ as
$\nparen{\Pr[X=\omega\emid Z=z]-\Pr[X=\omega]}
+\nparen{\Pr[X=\omega]-\Pr[W=\omega]}$ and square. The cross term
averages to $0$ over $z\sim Z$, since
$\E_{z\sim Z}\Pr[X=\omega\emid Z=z]=\Pr[X=\omega]$.
\end{proof}
Define the quadratic information of $X$ revealed by $Z$ as
\begin{equation*}
  \tsmuti{X}{Z}=\tsdiv{X\emid Z}{X}.
\end{equation*}
The same computation with $\Pr[W=\omega]$ replaced by $0$ gives
the chain identity
\begin{equation}
  \label{eq:quad-chain}
  \tsmuti{X}{Z}=\tsent{X}-\tsent{X\emid Z}.
\end{equation}
For a bit $B$ with $\E B=\mu$, this gives
$\tsent{B}=4\mu(1-\mu)$ and
$\tsmuti{B}{Z}=4\E\bigl[(\Pr[B=1\emid Z]-\mu)^2\bigr]$.
For instance, $\tsmuti{Y_1}{X}=\rho^2$, since
$\Pr[Y_1=1\emid X]=(1\pm\rho)/2$.

Let $X,Z,W$ be random variables on finite alphabets.
As before, conditioning means averaging:
\begin{equation*}
  \tsmuti{X}{W\emid Z}=\E_{z\sim Z}\tsmuti{X}{W\emid Z=z}
  =\tsent{X\emid Z}-\tsent{X\emid Z,W}.
\end{equation*}

\begin{lemma}
\label{lem:quad-properties}
\leavevmode
\begin{enumerate}
\item $0\le\tsmuti{X}{Z}\le\tsent{X}$, with
  $\tsmuti{X}{Z}=0$ iff $X$ and $Z$ are independent, and
  $\tsmuti{X}{Z}=\tsent{X}$ iff $X$ is a function of $Z$.
\item (Chain rule.) We have $\tsmuti{X}{W\emid Z}\ge0$ and
\begin{equation*}
  \tsmuti{X}{Z,W}=\tsmuti{X}{Z}+\tsmuti{X}{W\emid Z}.
\end{equation*}
\item (Data processing.) If $X-Z-W$ is a Markov chain, then
  $\tsmuti{X}{W}\le\tsmuti{X}{Z}$.
\end{enumerate}
\end{lemma}
\begin{proof}
(1) Nonnegativity is clear. The upper bound is
\eqref{eq:quad-chain} with $\tsent{X\emid Z}\ge0$, and
$\tsent{X\emid Z}=0$ iff $X$ is a function of $Z$.
(2) Telescope \eqref{eq:quad-chain}; the conditional term is the
average of $\tsmuti{X}{W\emid Z=z}$ over $z\sim Z$, hence
nonnegative.
(3) By the Markov property $\tsent{X\emid Z,W}=\tsent{X\emid Z}$,
so the chain rule in both orders gives
$\tsmuti{X}{Z}=\tsmuti{X}{Z,W}=\tsmuti{X}{W}+\tsmuti{X}{Z\emid W}
\ge\tsmuti{X}{W}$.
\end{proof}

Unlike Shannon information, $\tsmuti{\cdot}{\cdot}$ is not
symmetric. If $X$ is a uniform bit and $Z=(X,W)$ with $W$ an
independent uniform bit, then $\tsmuti{X}{Z}=\tsent{X}=1$ while
$\tsmuti{Z}{X}=\tsent{Z}-\tsent{Z\emid X}=\tfrac32-1=\tfrac12$.
This is why the notation uses a semicolon and why the order of
arguments matters in \autoref{sec:proof}: the first argument is
always the variable whose conditional law is compared with its marginal law.

\fi

\begin{lemma}
\label{lem:binary-divergence-comparison}
For $a\in[0,1]$ and $b\in[1/2,1)$,
\begin{equation*}
  \frac{2(a-b)^2}{\ln2}
  \le\kldivb{a}{b}
  \le\frac{\ln(1/b)}{(1-b)^2\ln2}(a-b)^2.
\end{equation*}
The lower bound holds for every $b\in(0,1)$; complementing
$a,b$ gives the corresponding upper bound for $b<1/2$.
\end{lemma}
\begin{proof}
The second derivative of the divergence in $a$ is
$1/(a(1-a)\ln2)\ge4/\ln2$, and its value and first derivative
vanish at $a=b$. This proves the lower bound.
For the upper bound, Taylor's integral formula gives
\begin{equation*}
  \frac{\kldivb{a}{b}}{(a-b)^2}
  =\frac1{\ln2}\int_0^1
  \frac{1-t}{(b+t(a-b))(1-b-t(a-b))}\,\mathrm dt.
\end{equation*}
The quotient extends continuously to $a=b$ and is convex in
$a$, because $1/(x(1-x))=1/x+1/(1-x)$ is convex.
It therefore attains its maximum at an endpoint.
For $b\ge1/2$, the value at $a=1$ is at least the value at
$a=0$: equivalently,
\begin{equation*}
  -b^2\ln b+(1-b)^2\ln(1-b)\ge0.
\end{equation*}
The expression vanishes at $b=1/2$ and $b=1$, and its second
derivative is $2\ln((1-b)/b)\le0$ on this interval.
The value at $a=1$ is the claimed coefficient.
\end{proof}

For bits, we write $\tsdivb{a}{b}=4(a-b)^2$ and
$\tsent{a}=4a(1-a)$.
The binary divergence has the expansion
\begin{equation}
  \label{eq:binary-divergence-series}
  \kldivb{\tfrac{1+u}{2}}{\tfrac12}
  =\sum_{k\ge1}\frac{u^{2k}}{2k(2k-1)\ln2},
  \qquad -1\le u\le1.
\end{equation}
Indeed, expanding its second derivative
$1/((1-u^2)\ln2)$ as a geometric series and integrating twice
proves the identity for $|u|<1$; continuity gives the endpoints.
The coefficients $c_k=1/(2k(2k-1)\ln2)$ are nonnegative and
satisfy $\sum_k c_k=1$ (set $u=1$). Hence,
\begin{equation*}
  \kldivb{\tfrac{1+u}2}{\tfrac12}=\sum_{k\ge1}c_ku^{2k}
  \le u^2=\tsdivb{\tfrac{1+u}2}{\tfrac12},
\end{equation*}
with equality iff $u\in\{-1,0,1\}$. Since
$\entb{a}=1-\kldivb{a}{1/2}$ and
$\tsent{a}=1-\tsdivb{a}{1/2}$, this gives
\begin{equation*}
  \tsent{a}\le\entb{a}\quad(0\le a\le1),
  \qquad
  \muti{B}{Z}\le\tsmuti{B}{Z}\quad\text{when $B$ is a uniform bit}.
\end{equation*}
In particular $\tsent{B}\le\ent{B}$ and
$\muti{Y_1}{X}=\sI\le\qI=\tsmuti{Y_1}{X}$,
the orderings used in \autoref{fig:q-square}. The second
inequality needs the uniform reference: for a nonconstant biased
bit $B$, taking $Z=B$ gives
$\muti{B}{B}=\ent{B}>\tsent{B}=\tsmuti{B}{B}$.

\begin{proof}[Proof of \autoref{lem:tsallis-properties}]
Let us write $\mu_Z=\Pr[C=1\emid Z]$ and $\mu=\E C$.
Expanding the definition gives
\begin{equation*}
  \tsmuti{C}{Z}=4\E\bigl[(\mu_Z-\mu)^2\bigr].
\end{equation*}
This proves nonnegativity, while the upper bound by $\tsent{C}$
follows from $\tsent{C\emid Z}\ge0$. Both statements may be
averaged over $U$.
We already have $\tsent{C}\le\ent{C}$. The lower bound of
\autoref{lem:binary-divergence-comparison} with reference mean
$1/2$ also gives
\begin{equation*}
  1-\ent{C}=\kldivb{\mu}{1/2}
  \ge\frac{1-\tsent{C}}{2\ln2}
  \ge\frac{1-\tsent{C}}2.
\end{equation*}
Averaging over $U$ completes part~\ref{item:tsallis-bounds}.

The chain rule follows by telescoping the entropy differences.
If $C-Z-W$ is a Markov chain, then
$\tsent{C\emid Z,W}=\tsent{C\emid Z}$, so the chain rule
also gives the data processing inequality in
part~\ref{item:tsallis-chain}.

The lower bound of \autoref{lem:binary-divergence-comparison},
averaged over the conditional mean $\mu_Z$, gives
\begin{equation*}
  \muti{C}{Z}\ge\frac{2\E[(\mu_Z-\mu)^2]}{\ln2}
  =\frac{\tsmuti{C}{Z}}{2\ln2}.
\end{equation*}
The upper comparison for unbiased $C$ was proved above using
\eqref{eq:binary-divergence-series}. Averaging both comparisons
over $U$ gives part~\ref{item:tsallis-shannon}; the upper one
requires $C$ to remain unbiased given $U$. If a conditional law
of $C$ is constant, both informations are zero.

For two bits, evaluating the two conditional means gives
\begin{equation*}
  \tsent{C}\tsmuti{B}{C}
  =16\operatorname{Cov}(B,C)^2
  =\tsent{B}\tsmuti{C}{B}.
\end{equation*}
If either bit is constant, all three terms vanish.
When $B$ is unbiased, the definition also gives
\begin{equation*}
  \tsmuti{C}{B}
  =\bigl(\Pr[C=1\emid B=1]-\Pr[C=1\emid B=0]\bigr)^2
  =\bigl(2\Pr[B=C]-1\bigr)^2.
\end{equation*}
This proves part~\ref{item:tsallis-bits}.

For part~\ref{item:tsallis-conditioning}, let us write
$V_u=\tsent{C\emid U=u}$ and
$\sigma_u=\operatorname{Cov}(B,C\emid U=u)$.
The identity for two bits, applied conditionally, and
Cauchy--Schwarz give
\begin{align*}
  \tsmuti{B}{C\emid U}\,\tsent{C\emid U}
  &=16\left(\E_{u\sim U}\frac{\sigma_u^2}{V_u}\right)
  \left(\E_{u\sim U}V_u\right)\\
  &\ge16\left(\E_{u\sim U}\sigma_u\right)^2.
\end{align*}
Terms with $V_u=0$ have $\sigma_u=0$ and contribute zero.
Since $B$ is independent of $U$, its conditional mean is constant,
and hence $\E_{u\sim U}\sigma_u=\operatorname{Cov}(B,C)$.
The last display is therefore at least
$\tsent{B}\tsmuti{C}{B}$, as required.

Finally, applying the Shannon comparisons to the unbiased noise
bits gives $\qI/(2\ln2)\le\sI\le\qI$. Since $\ln2<1$,
we have $\qI\le2\sI$, proving part~\ref{item:tsallis-noise}.
\end{proof}

\begin{proof}[Proof of \autoref{lem:biased-comparison}]
Complementing $C$ if necessary lets us assume $\mu\ge1/2$.
Averaging the upper bound of
\autoref{lem:binary-divergence-comparison} over the conditional
mean $\Pr[C=1\mid U]$ gives the first inequality.
For the second,
\begin{equation*}
  \tsent{C}\,c(\mu)
  =\frac{\mu}{1-\mu}\ln\left(1+\frac{1-\mu}{\mu}\right)
  \le1.
\end{equation*}
using $\ln(1+x)\le x$.
\end{proof}
